\section{Pricing and Return on Investment}
\label{sec:pricing}

\subsection{How Pricing Works}

\sys{} uses a transparent usage-based pricing model.
Every session costs exactly what the resources consumed cost — no subscriptions
with unused capacity, no surprise overruns.

The cost of a session is determined by two factors:

\begin{enumerate}
  \item \textbf{Text generated.} Longer documents with more sections cost more
        than shorter ones.
        The cost per thousand words of generated text is approximately \$0.30--0.50
        at current API rates.
  \item \textbf{Figures generated.} Each figure requires compilation and potentially
        multiple repair attempts.
        The cost per figure is approximately \$0.08--0.12.
\end{enumerate}

A typical 20-page academic document with 30 references and 4 figures costs
between \$2.50 and \$6.00 in total session cost.

\begin{table}[h!]
\centering
\caption{Typical session costs by document type.}
\label{tab:pricing}
\begin{tabular}{lrrr}
\toprule
Document type & Length & Figures & Estimated cost \\
\midrule
Short report / abstract     & 3--5 pages  & 0--1 & \$0.50--1.50 \\
Conference paper draft       & 8--12 pages & 2--4 & \$1.50--4.00 \\
Journal article draft        & 15--25 pages & 4--8 & \$3.00--8.00 \\
Literature review chapter    & 10--20 pages & 2--5 & \$2.00--6.00 \\
Full thesis chapter          & 20--40 pages & 5--12 & \$5.00--15.00 \\
\bottomrule
\end{tabular}
\end{table}

\subsection{Institutional Pricing}

For universities and organisations purchasing access for multiple users,
\sys{} offers pre-loaded credit packages with volume discounts.
An organisation that pre-loads \$500 in credits receives access for the equivalent
of 80--200 full document sessions, depending on document length and figure count.

Institutional accounts additionally receive:
\begin{itemize}
  \item A shared billing dashboard visible to department administrators.
  \item Per-user budget caps (so individual researchers cannot accidentally
        exceed their allocated quota).
  \item Session history and audit logs for compliance purposes.
  \item Priority processing (sessions run on dedicated capacity, not shared queues).
\end{itemize}

\subsection{Return on Investment}

The ROI calculation for \sys{} is straightforward.
The cost of a human researcher's time doing equivalent work is the baseline.

\begin{table}[h!]
\centering
\caption{ROI calculation: \sys{} vs.\ researcher time for a 20-page document.}
\label{tab:roi}
\begin{tabular}{lrr}
\toprule
Cost component & Manual workflow & \sys{} \\
\midrule
Researcher time (hours)              & 20--40 hrs & 1--2 hrs (revision) \\
Researcher time cost (\$50/hr avg.)  & \$1,000--2,000 & \$50--100 \\
Platform session cost                & \$0 & \$3--6 \\
External writing service (if used)   & \$200--800 & \$0 \\
\midrule
\textbf{Total cost per document}     & \textbf{\$1,000--2,800} & \textbf{\$53--106} \\
\midrule
\textbf{Cost reduction}              & — & \textbf{90--96\%} \\
\bottomrule
\end{tabular}
\end{table}

The table above values researcher time at \$50/hr — a conservative figure for a
postdoctoral researcher or junior R\&D employee in the CIS region.
For senior researchers or those in Western institutions, the figure is 2--5$\times$ higher,
making the ROI proportionally larger.

\begin{statbox}[title=Break-Even Analysis for a University Department]
Assumptions: 20 active researchers, each producing 4 documents per year,
saving 15 hours per document at \$50/hr researcher cost.\\[0.4em]
\textbf{Annual time saved:} 20 × 4 × 15 = 1,200 researcher-hours\\
\textbf{Annual cost saved:} 1,200 × \$50 = \$60,000\\
\textbf{Annual platform cost:} 80 sessions × \$5 avg = \$400\\[0.4em]
\textbf{Return on platform investment: 150$\times$}
\end{statbox}

\subsection{What \sys{} Does Not Replace}

It is important to be specific about what \sys{} does and does not do, so that
purchasing expectations are calibrated correctly.

\sys{} \textbf{does} produce:
a verified-reference first draft, compilable figures, correctly formatted bibliography,
and a clean PDF.

\sys{} \textbf{does not} replace:
the researcher's own experimental results and analysis,
the researcher's interpretation and scientific judgment,
the peer review process,
or the final revision and submission workflow.

The appropriate framing is: \sys{} handles the parts of academic writing that are
formulaic, time-consuming, and error-prone.
The researcher handles the parts that require expertise, originality, and judgment.
Together, they produce a document that is both faster to create and more reliable
than either could produce alone.

\begin{figure}[h!]
\centering
\begin{tikzpicture}
\begin{axis}[
  PL,
  ybar, bar width=28pt,
  ylabel={Cost per 20-page document (USD)},
  title={Cost comparison: document production methods},
  symbolic x coords={
    {Human service\\(CIS)},
    {Human service\\(Western)},
    {Manual AI\\workflow},
    {\sys{}}
  },
  xtick=data,
  x tick label style={align=center, font=\small, text width=2.8cm},
  ymin=0, ymax=1300,
  ytick={0,200,400,600,800,1000,1200},
  yticklabel={\$\pgfmathprintnumber[fixed, precision=0]{\tick}},
  width=13cm, height=7cm,
  nodes near coords,
  nodes near coords style={font=\scriptsize\bfseries, anchor=south},
  every node near coord/.append style={
    /pgf/number format/.cd, fixed, precision=0,
    /tikz/.cd
  },
]
\addplot[
  draw=C4, fill=F4, fill opacity=0.9,
  nodes near coords={\$\pgfmathprintnumber[fixed,precision=0]{\pgfplotspointmeta}},
] coordinates {
  ({Human service\\(CIS)},500)
  ({Human service\\(Western)},1200)
  ({Manual AI\\workflow},80)
  ({\sys{}},5)
};
\end{axis}
\end{tikzpicture}
\caption{Total cost per 20-page academic document by production method,
         including labour time valued at local market rates.
         ``Manual AI workflow'' includes the researcher time spent prompting,
         verifying references manually, creating figures, and fixing compilation errors.
         \sys{} costs reflect the platform session fee only; light revision time
         (1--2 hours) is excluded for comparability.}
\label{fig:roi}
\end{figure}

